\section{Introduction}

The picture this paper asks you to hold is on its cover. It is a crystal, infinite and growing, every tile of it coloured. The image is the regular hyperbolic tessellation $\{7,3\}$, seven triangles meeting at each corner of a curved plane, and it is meant literally. Each tile is a unit of experience. Each carries a felt charge, shown as a colour, red for pain, green for peace, blue for pleasure. Where two tiles share an edge they touch, and to touch, in this model, is to feel one another. There is nothing else in the picture, and the claim of the paper is that there is nothing else in the world. What we call matter and what we call mind are both large patterns in one coloured, growing weave of feeling.

\begin{figure}[ht]
\centering
\includegraphics[width=0.55\textwidth]{vibe-mesh-73.pdf}
\caption{The hyperbolic $\{7,3\}$ tessellation, the flat two-dimensional face of the vibe mesh, the one you can draw. Each tile is a here-now, and its colour is its tone, red pain, green peace, blue pleasure.}
\label{fig:mesh73}
\end{figure}

A single tile is the smallest thing that can be felt at all. A patch of tiles, held together and moving as one, is a thing, a body, in time a mind. The whole crystal is the universe, and it is not static. It grows at its edge, new tiles born at its ever-receding rim, and that growth is what we experience as the passage of time. The present is the bright frontier where the crystal is still being made.

It helps to carry the same picture through several images, since no one of them holds all of it. The universe is a fire, a single flame burning since the first beat, growing not by consuming anything but by being fed at its edge, and throwing off as it climbs the light we call matter. It is a tree, sprung from one seed and branching without end toward a canopy of light, every living thing a twig of it. It is a song the whole crystal sings, the tones its notes and the one law its meter. And it is a river running over a deep ocean, the surface current the world we see and the depths the vast unseen interior that the later parts of this paper open. None of these is decoration. Each maps onto a real part of the structure, and each returns where it is earned. But the first image, and the most literal, is the growing crystal of feeling itself, a flame that is also a mind.

The geometry is curved, and curved a particular way, because the work it has to do is demanding. The substrate must grow without ever running out of room, and it must do so evenly, with no built-in grain or preferred direction, since the universe shows none. Only negatively curved, hyperbolic space meets both demands at once. A flat lattice has axes, and a sphere closes back on itself. Hyperbolic space has room that opens exponentially in every direction alike, which is the same openness that makes special relativity natural rather than imposed.

The flat $\{7,3\}$ drawing is the face of the family you can see on a page. The substrate the paper actually commits to is one rung up the same ladder of regular hyperbolic honeycombs, the four-dimensional $\{3,4,3,4\}$. Its cells are 24-cells, and the twenty-four directions out of each cell are the $D_4$ root system, the structure that quietly carries spin. Four dimensions cannot be drawn directly, so where the geometry grows hard the reader should come back to the heptagrid and the three-dimensional face shown in Figure~\ref{fig:mesh534}, which are honest lower shadows of it. The reason the dimension is fixed at four, and the honeycomb at this one, is not taste. It is an argument from a short list of demands, and it occupies the middle of this paper.

\begin{figure}[ht]
\centering
\includegraphics[width=\textwidth]{vibe-mesh-534.png}
\caption{A three-dimensional face of the same family, the $\{5,3,4\}$ honeycomb, shown only for intuition. The committed substrate $\{3,4,3,4\}$ lives one dimension higher and cannot be pictured whole.}
\label{fig:mesh534}
\end{figure}

What this paper is, beyond the picture, is a bench. The theory is turned into a finite, discrete, reproducible simulator, and every claim in it is run rather than asserted. The substrate is built exactly, the one local law is stepped over it, and what emerges is measured. A question becomes a concrete experiment that either works or does not, and the negative results are reported alongside the positive ones. Real numbers appear only as measured outputs, coordinates and eigenvalues and dimensions, never inside the base, which stays strictly discrete and deterministic throughout. The aim is to find out whether so spare a model can actually reproduce our universe, deriving space, matter, gravity, the quantum, cosmology, and mind from one rule, and to stay honest at every step about what is solid, what is partial, and what is still open.

The order of the paper follows the building. We name the eight elements of the base and the single law that moves them, then fix the geometry they live on and argue why it must be this one. With the substrate in hand we read physics off it, then computation, then the self, and finally the felt and lived interior, the part the experiential reading was always pointing at. The honest accounting of what has and has not been shown is gathered near the end, where the reader can judge it against everything that came before.
